\subsection{Preliminaries on automorphisms of rings of formal power series}

LEMMA 2. Let \(f(\mathbf{X}_1, \mathbf{X}_2, \ldots, \mathbf{X}_n)\) be a formal power series \(\neq 0\) with coefficients in a ring E. There exist integers \(u(i) \geqslant 1\) ( \(1 \leqslant i \leqslant n - 1\) ) such that

\[
f (\mathrm{T} ^ {u (1)}, \dots , \mathrm{T} ^ {u (n - 1)}, \mathrm{T}) \neq 0.
\]

Suppose that integers \(u(i) \geqslant 1\) ( \(1 \leqslant i \leqslant k - 1\) ) are determined such that \(f(\mathbf{X}_n^{u(1)}, \ldots, \mathbf{X}_n^{u(k-1)}, \mathbf{X}_k, \ldots, \mathbf{X}_n) \neq 0\) . We shall determine an integer \(u(k) \geqslant 1\) such that

\[
f \left(\mathrm{X} _ {n} ^ {u (1)}, \dots , \mathrm{X} _ {n} ^ {u (k - 1)}, \mathrm{X} _ {n} ^ {u (k)}, \mathrm{X} _ {k + 1}, \dots , \mathrm{X} _ {n}\right) \neq 0.
\]

The lemma will then be proved by induction.

Observe that the series \(f(\mathbf{X}_n^{u(1)}, \ldots, \mathbf{X}_n^{u(k-1)}, \mathbf{X}_k, \ldots, \mathbf{X}_n)\) can be considered as a series in \(\mathbf{X}_k\) and \(\mathbf{X}_n\) with coefficients in \(\mathbf{E}[[\mathbf{X}_{k+1}, \ldots, \mathbf{X}_{n-1}]]\) . Thus we see that it suffices to establish the lemma for \(n = 2\) .

Therefore let

\[
f = \sum_ {i, j} e _ {i j} \mathbf {X} ^ {i} \mathbf {Y} ^ {j} \in \mathrm{E} [ [ \mathbf {X}, \mathbf {Y} ] ]
\]

where \(f \neq 0\) . Let \(G \subset N \times N\) be the non-empty set of ordered pairs \((i, j)\) such that \(e_{ij} \neq 0\) . Let \(N \times N\) be given the lexicographical ordering. Let \((c, d)\) be the least element of \(G\) . Choose an integer \(p > d\) . In the expansion of

\[
f (\mathbf {T} ^ {p}, \mathbf {T}) = \sum_ {(i, j) \in \mathbb {G}} e _ {i j} \mathbf {T} ^ {i p + j}
\]

we look for the terms of degree \(cp + d\) . If \(ip + j = cp + d\) , \(i \geqslant c + 1\) is impossible, for this would give

\[
i p + j \geqslant (c + 1) p + j \geqslant (c + 1) p > c p + d;
\]

nor is \(i < c\) possible, for \((c, d)\) is the least element of \(G\) ; therefore \(i = c\) and then \(j = d\) . The term of degree \(cp + d\) in \(f(T^p, T)\) is therefore \(e_{cd}T^{cp + d}\) . Since \(e_{cd} \neq 0, f(T^p, T) \neq 0\) . Whence the lemma.

In the ring \(\operatorname{E}[[\mathbf{X}_1, \ldots, \mathbf{X}_n]]\) , let \(a\) be the ideal of formal power series without constant term. If \(w_1, \ldots, w_n\) are elements of \(a\) , recall that the mapping \(f(\mathbf{X}_1, \ldots, \mathbf{X}_n) \mapsto f(w_1, \ldots, w_n)\) is the unique endomorphism \(s\) of the ring \(\operatorname{E}[[\mathbf{X}_1, \ldots, \mathbf{X}_n]]\) such that \(s(\mathbf{X}_i) = w_i\) for \(1 \leqslant i \leqslant n\) (Chapter III, § 4, no. 5, Proposition 6).

We take \(w_1 = \mathbf{X}_1 + \mathbf{X}_n^{u(1)}, \ldots, w_{n-1} = \mathbf{X}_{n-1} + \mathbf{X}_n^{u(n-1)}, w_n = \mathbf{X}_n\) , where the \(u(i)\) are integers \(\geqslant 1\) . Let \(s'\) be the endomorphism of \(\operatorname{E}[[\mathbf{X}_1, \ldots, \mathbf{X}_n]]\) which maps \(\mathbf{X}_1\) to \(\mathbf{X}_1 - \mathbf{X}_n^{u(1)}, \ldots, \mathbf{X}_{n-1}\) to \(\mathbf{X}_{n-1} - \mathbf{X}_n^{u(n-1)}\) and \(\mathbf{X}_n\) to \(\mathbf{X}_n\) . Then \(s'(s(\mathbf{X}_i)) = \mathbf{X}_i\) for \(1 \leqslant i \leqslant n\) and hence \(s' \circ s\) is the identity automorphism; similarly for \(s \circ s'\) . Hence \(s\) is an automorphism.

LEMMA 3. Let \(f\) be a non-zero element of \(\operatorname{E}[[\mathbf{X}_1, \ldots, \mathbf{X}_n]]\) . There exist integers \(u(i) \geqslant 1 (1 \leqslant i \leqslant n - 1)\) such that the automorphism \(s\) of \(\operatorname{E}\) defined by

\[
s (\mathbf {X} _ {i}) = \mathbf {X} _ {i} + \mathbf {X} _ {n} ^ {u (i)} \quad (1 \leqslant i \leqslant n - 1)
\]

and \(s(\mathbf{X}_n) = \mathbf{X}_n\) maps \(f\) to an element \(g\) such that \(g(0, \ldots, 0, \mathbf{X}_n) \neq 0\) .

\(g(0,\ldots,0,\mathbf{X}_n) = f(\mathbf{X}_n^{u(1)},\ldots ,\mathbf{X}_n^{u(n - 1)},\mathbf{X}_n)\) . Lemma 3 is therefore a consequence of Lemma 2.

